\section{Remote providers, local ingestion, and search}
\label{sec:appendix-remote-ingestion}
\label{sec:app-remote-ingestion}

Use the external relations to discover records and retrieve a bounded set of
sequences. Store the sequences in a local BioSeqDB table before running repeated
searches. The local copy gives the application a fixed search target. Separate
requests to a provider do not share a PostgreSQL snapshot.
The external relations are read-only. They reject \texttt{INSERT},
\texttt{UPDATE}, \texttt{DELETE}, \texttt{COPY FROM}, and
\texttt{CREATE INDEX}. Use a local table as the write destination.

The examples below use a scratch schema or an otherwise empty working schema.
Run them in order, using a role that can create local tables. Provider examples
require network access from the database server. The local rehearsal requires
no network access. Identifiers and response metadata belong to their providers.
The examples do not promise that a live service will return a particular row.

\subsection{Check installation and permissions}
\label{app-remote-permissions}

\begin{enumerate}
  \item Install BioSeqDB with external-provider support enabled. The ordinary
        local sequence extension alone does not establish this capability.
  \item Check the schema, table access method, and six installed relations.
        This query reads PostgreSQL catalogs and makes no remote request.
\end{enumerate}

\par\noindent\begin{minipage}{\linewidth}
\begin{lstlisting}[style=bioseq-appendix-sql]
SELECT
  EXISTS (
    SELECT 1 FROM pg_namespace
    WHERE nspname = 'bioseq_external'
  ) AS schema_ready,
  EXISTS (
    SELECT 1 FROM pg_am WHERE amname = 'bioseq_external'
  ) AS table_am_ready,
  (
    SELECT count(*)
    FROM pg_class AS c
    JOIN pg_namespace AS n ON n.oid = c.relnamespace
    WHERE n.nspname = 'bioseq_external'
      AND c.relname IN (
        'taxonomy', 'catalog', 'record_genes',
        'record_locations', 'sequences', 'sequence_chunks'
      )
  ) AS relation_count;
\end{lstlisting}
\end{minipage}

\par\noindent\begin{minipage}{\linewidth}
\begin{lstlisting}[style=bioseq-appendix-sql]
SELECT c.relname, a.amname
FROM pg_class AS c
JOIN pg_namespace AS n ON n.oid = c.relnamespace
LEFT JOIN pg_am AS a ON a.oid = c.relam
WHERE n.nspname = 'bioseq_external'
ORDER BY c.relname;
\end{lstlisting}
\end{minipage}

The first query should report two true values and a count of six for the
documented installation. Relation existence does not establish connectivity or
permission to read it. The schema and its tables are restricted by default.
An administrator can grant an existing application role access as follows.
Replace \texttt{sequence\_importer} with that role's name.

\par\noindent\begin{minipage}{\linewidth}
\begin{lstlisting}[style=bioseq-appendix-sql]
GRANT USAGE ON SCHEMA bioseq_external TO sequence_importer;

GRANT SELECT ON
  bioseq_external.taxonomy,
  bioseq_external.catalog,
  bioseq_external.record_genes,
  bioseq_external.record_locations,
  bioseq_external.sequences,
  bioseq_external.sequence_chunks
TO sequence_importer;
\end{lstlisting}
\end{minipage}

Grant only the relations the application needs. Permission to select an
external relation permits the associated server-side network requests.
The current SQL API has no \texttt{taxon\_matches} or
\texttt{native\_query} marker function to grant or call.

\subsection{Discover taxonomy and catalog metadata}
\label{app-remote-discovery}

Begin with an explicit provider and a narrow lookup. A name lookup can return
provider taxonomy identifiers. An exact numeric identifier avoids ambiguity
between names. Neither query is a descendant-taxonomy expansion.

\par\noindent\begin{minipage}{\linewidth}
\begin{lstlisting}[style=bioseq-appendix-sql]
SELECT provider, taxon_id, provider_taxon_id,
       scientific_name, common_name, rank
FROM bioseq_external.taxonomy
WHERE provider = 'ncbi'
  AND scientific_name = 'Homo sapiens'
LIMIT 5;

SELECT provider, taxon_id, scientific_name, rank
FROM bioseq_external.taxonomy
WHERE provider = 'ncbi' AND taxon_id = 9606
LIMIT 5;
\end{lstlisting}
\end{minipage}

Use \texttt{catalog} when you want descriptions, lengths, or retrieval
capabilities without sequence data. UniProt supplies protein metadata here.
The NCBI catalog route supplies genome assembly reports, so use an assembly
accession beginning with \texttt{GCF\_} or \texttt{GCA\_} for that route.
NCBI Datasets distinguishes these reports from sequence download
packages~\autocite{ncbi_datasets_v2_api}.

\par\noindent\begin{minipage}{\linewidth}
\begin{lstlisting}[style=bioseq-appendix-sql]
SELECT provider, accession, molecule_kind, taxon_id,
       scientific_name, title, sequence_length, retrievable
FROM bioseq_external.catalog
WHERE provider = 'uniprot' AND accession = 'P04637';

SELECT provider, accession, title, assembly_level, retrievable
FROM bioseq_external.catalog
WHERE provider = 'ncbi' AND accession = 'GCF_000001405.40';
\end{lstlisting}
\end{minipage}

Small accession arrays can express several lookups. Current remote predicate
extraction requires constants in the plan. For a prepared remote query, force
a custom plan for each execution so PostgreSQL can substitute the values.
Do not rely on a generic plan to bound the provider request.

\par\noindent\begin{minipage}{\linewidth}
\begin{lstlisting}[style=bioseq-appendix-sql]
BEGIN;
SET LOCAL plan_cache_mode = force_custom_plan;

PREPARE external_metadata(text[]) AS
SELECT provider, accession, title, sequence_length
FROM bioseq_external.catalog
WHERE provider IN ('ncbi', 'uniprot', 'ebi')
  AND taxon_id = 9606
  AND accession = ANY ($1)
ORDER BY provider, accession;

EXECUTE external_metadata(
  ARRAY['P04637', 'GCF_000001405.40']
);
DEALLOCATE external_metadata;
COMMIT;
\end{lstlisting}
\end{minipage}

This mixed-provider example sends the accession list to each selected
provider. The identifiers need not be valid in every provider. Use separate
provider-specific queries when you need to verify each source independently.
The scan expands at most 64 combinations of non-provider predicates per
provider. Arrays cause separate lookups, not one guaranteed batched HTTP
request. Combinations beyond the cap are omitted with a warning.

\subsection{Choose an operation the provider supports}

\begin{description}
  \item[NCBI] Taxonomy and genome assembly discovery, gene locations,
        supported nucleotide retrieval, and sequence chunks.
  \item[UCSC] Assembly discovery and nucleotide regions identified by
        assembly database, chromosome, and coordinates.
  \item[UniProt] Protein catalog metadata and record-to-gene expansion.
        Protein sequence retrieval is not implemented by this SQL route.
  \item[EBI] Discovery metadata. This route does not retrieve sequence
        characters.
\end{description}

An unsupported sequence operation can raise SQLSTATE \texttt{0A000}.
When several lookups are combined, a successful statement can include failed
provider requests, even when later SQL filters leave zero rows. Statement
success therefore does not prove that every provider completed. Catalog
discovery also has provider page and result limits. \texttt{LIMIT} is not an
exhaustive paging
interface. Prefer small, explicit requests and check their coverage locally.

\subsection{Expand genes and genomic locations}

The child relations expose repeated metadata as rows. Their child indexes
are zero-based. Their count columns describe the available children for the
record. The following examples inspect UniProt gene annotations and locations
for NCBI Gene ID 7157.

\par\noindent\begin{minipage}{\linewidth}
\begin{lstlisting}[style=bioseq-appendix-sql]
SELECT accession, gene_index, gene_count,
       gene_namespace, gene_id, gene_symbol, gene_name
FROM bioseq_external.record_genes
WHERE provider = 'uniprot' AND accession = 'P04637'
ORDER BY gene_index;

SELECT accession, location_index, location_count,
       assembly_accession, assembly_database, sequence_alias,
       start0, end0, strand
FROM bioseq_external.record_locations
WHERE provider = 'ncbi' AND accession = '7157'
  AND start0 < end0
ORDER BY location_index
LIMIT 2;
\end{lstlisting}
\end{minipage}

Location intervals use zero-based, half-open coordinates:
\([\texttt{start0},\texttt{end0})\). Their length is
\texttt{end0 - start0}. Retain the assembly and strand before using a location
to request sequence. A gene identifier alone is not a chromosome coordinate.

\subsection{Inspect the plan before requesting data}

Plain \texttt{EXPLAIN} shows the external scan, its entity, and the selected
provider without executing the request. Adding \texttt{ANALYZE} executes it.

\par\noindent\begin{minipage}{\linewidth}
\begin{lstlisting}[style=bioseq-appendix-sql]
EXPLAIN (COSTS OFF)
SELECT taxon_id, scientific_name
FROM bioseq_external.taxonomy
WHERE provider = 'ncbi' AND taxon_id = 9606;

EXPLAIN (ANALYZE, COSTS OFF)
SELECT taxon_id, scientific_name
FROM bioseq_external.taxonomy
WHERE provider = 'ncbi' AND taxon_id = 9606;
\end{lstlisting}
\end{minipage}

Inspect the first plan before choosing to run the second. A sequence relation
retrieves sequence data even when the projection only asks for an accession
or a count. Select \texttt{catalog} for a metadata-only operation.

\subsection{Retrieve a bounded nucleotide region}

The following query asks UCSC for the first 100 bases of \texttt{hg38 chrM}.
Supply both interval bounds. These are the same half-open coordinates used
by UCSC's sequence endpoint~\autocite{ucsc_rest_api}.

\par\noindent\begin{minipage}{\linewidth}
\begin{lstlisting}[style=bioseq-appendix-sql]
SELECT provider, accession, assembly_database, chromosome,
       start0, end0, bioseq.len(sequence) AS bases
FROM bioseq_external.sequences
WHERE provider = 'ucsc'
  AND assembly_database = 'hg38'
  AND chromosome = 'chrM'
  AND start0 = 0 AND end0 = 100;
\end{lstlisting}
\end{minipage}

The returned \texttt{bioseq.seq} value is unbound: it contains sequence
characters before a destination alphabet and packed representation are
selected. A cast or assignment to a bound local sequence column supplies
those choices. The remote relation does not establish \texttt{dna5/64}
storage merely because its value has the SQL type \texttt{bioseq.seq}.

For a whole NCBI nucleotide record, a small local DNA5 table can accept the
returned value by direct assignment. The destination column chooses the
alphabet and storage unit, then validates and encodes the characters.

\par\noindent\begin{minipage}{\linewidth}
\begin{lstlisting}[style=bioseq-appendix-sql]
CREATE TABLE imported_ncbi_records (
  accession text NOT NULL,
  sequence bioseq.seq('dna5', 64) NOT NULL
) USING bioseq;

INSERT INTO imported_ncbi_records (accession, sequence)
SELECT accession, sequence
FROM bioseq_external.sequences
WHERE provider = 'ncbi' AND accession = 'NC_012920.1';

SELECT accession, bioseq.len(sequence) AS bases
FROM imported_ncbi_records;
\end{lstlisting}
\end{minipage}

\subsection{Read a nucleotide record as ordered chunks}
\label{app-remote-chunks}

For a known NCBI nucleotide accession, the direct route uses EFetch FASTA
retrieval. EFetch accepts accession-version identifiers directly. No search
request is needed to resolve a known identifier~\autocite{ncbiEutilities}.
The same custom-plan rule applies to a prepared chunk query.

\par\noindent\begin{minipage}{\linewidth}
\begin{lstlisting}[style=bioseq-appendix-sql]
BEGIN;
SET LOCAL plan_cache_mode = force_custom_plan;

PREPARE external_chunks(text) AS
SELECT original_identifier, chunk_index, base_offset,
       bioseq.len(sequence_chunk) AS bases, end_of_record
FROM bioseq_external.sequence_chunks
WHERE provider = 'ncbi' AND accession = $1
ORDER BY chunk_index;

EXECUTE external_chunks('NC_012920.1');
DEALLOCATE external_chunks;
COMMIT;
\end{lstlisting}
\end{minipage}

Nonfinal chunks on this route contain 1,048,572 bases. The final chunk may be
shorter. \texttt{base\_offset} is relative to the current sequence record,
and \texttt{end\_of\_record} marks its last chunk. A short mitochondrial
record is sufficient to check the interface even though it needs only one
chunk. Keep \texttt{ORDER BY} when consuming ordered data: the external plan
does not advertise this ordering, so PostgreSQL may add a local sort.
Streaming the provider response therefore does not guarantee that the
complete SQL plan uses bounded memory or emits its first row immediately.

An exact NCBI assembly accession selects the Datasets archive route instead:

\par\noindent\begin{minipage}{\linewidth}
\begin{lstlisting}[style=bioseq-appendix-sql]
SELECT original_identifier, chunk_index, base_offset,
       bioseq.len(sequence_chunk) AS bases, end_of_record
FROM bioseq_external.sequence_chunks
WHERE provider = 'ncbi' AND accession = 'GCF_000005845.2'
ORDER BY chunk_index;
\end{lstlisting}
\end{minipage}

An assembly can contain several sequence records. Keep
\texttt{original\_identifier} to distinguish them: the accession identifies
the requested assembly, while offsets restart for each record. The current
archive route buffers a ZIP download of at most 64 MiB and allows at most
512 MiB decoded data and 4096 entries. It buffers each selected member and
checks its MD5 digest before emitting that member's sequence rows. It is not
a general large-archive importer.

\subsection{Create and rehearse the local staging workflow}
\label{app-remote-stage}

\begin{enumerate}
  \item Create a local table using the \texttt{bioseq} access method. This
        example accepts only DNA4 and stores region coordinates beside it.
  \item Insert a known local fixture and search it before making a live
        import. This separates installation or search problems from network
        problems.
\end{enumerate}

\par\noindent\begin{minipage}{\linewidth}
\begin{lstlisting}[style=bioseq-appendix-sql]
CREATE TABLE imported_provider_targets (
  provider text NOT NULL,
  accession text NOT NULL,
  start0 bigint NOT NULL,
  end0 bigint NOT NULL,
  sequence bioseq.seq('dna4', 64) NOT NULL,
  imported_at timestamptz NOT NULL DEFAULT clock_timestamp(),
  PRIMARY KEY (provider, accession, start0, end0)
) USING bioseq;

SET bioseq.index.insert_maintenance.enabled = on;
SET bioseq.vectorization.enabled = on;

INSERT INTO imported_provider_targets
  (provider, accession, start0, end0, sequence)
VALUES (
  'local-check', 'region-1', 1000, 1058,
  'AAAATAGTCACTGTAGCCTGATCGTAGCTAGTCAACGTACGTACGTACGTACGTAAAA'
    ::bioseq.seq('dna4', 64)
);
\end{lstlisting}
\end{minipage}

The maintenance switch permits each index's configured insert-update policy.
The vectorization switch permits vectorized SQL predicate and aggregate work.
It does not choose the minimizer or chaining algorithm.

The fixture is 58 bases long. Its 30-base query occurs at local interval
\([4,34)\), which maps to provider interval \([1004,1034)\). An exact
alignment with CIGAR \texttt{30M} is the expected check result. The helper
calls below remain plain projections inside the materialized CTE. Coordinate
arithmetic is performed by the outer query.

\par\noindent\begin{minipage}{\linewidth}
\begin{lstlisting}[style=bioseq-appendix-sql]
WITH local_matches AS MATERIALIZED (
  SELECT provider, accession, start0,
         bioseq.target_start(accession) AS match_start_in_region,
         bioseq.target_end(accession) AS match_end_in_region,
         bioseq.cigar(accession) AS cigar
  FROM imported_provider_targets
  WHERE public.mmsearch(
    sequence,
    'TAGTCACTGTAGCCTGATCGTAGCTAGTCA'::bioseq.seq('dna4', 64),
    aligner_options(best_n => 3)
  )
)
SELECT provider, accession, match_start_in_region,
       match_end_in_region, start0 + match_start_in_region
         AS provider_start0,
       start0 + match_end_in_region AS provider_end0, cigar
FROM local_matches
ORDER BY provider, accession, provider_start0;

DELETE FROM imported_provider_targets
WHERE provider = 'local-check';
\end{lstlisting}
\end{minipage}

\subsection{Import and search a remote region locally}

Use an explicit character policy when storing DNA4. The following query
excludes rows containing any symbol outside \texttt{A,C,G,T}. It does not
silently replace ambiguous bases. The direct cast binds an unbound remote
value to the local destination type. For N-bearing data, choose an
appropriate alphabet, such as DNA5, for both storage and subsequent search.

\par\noindent\begin{minipage}{\linewidth}
\begin{lstlisting}[style=bioseq-appendix-sql]
INSERT INTO imported_provider_targets
  (provider, accession, start0, end0, sequence)
SELECT provider, accession, start0, end0,
       sequence::bioseq.seq('dna4', 64)
FROM bioseq_external.sequences
WHERE provider = 'ucsc'
  AND assembly_database = 'hg38'
  AND chromosome = 'chrM'
  AND start0 = 0 AND end0 = 100
  AND sequence::text ~ '^[ACGT]*$';

SELECT provider, accession, start0, end0,
       bioseq.len(sequence) AS bases
FROM imported_provider_targets
ORDER BY provider, accession, start0;
\end{lstlisting}
\end{minipage}

Check the local row count and length, including whether the character policy
excluded a row. This ordinary \texttt{INSERT SELECT} uses PostgreSQL's
executor and the destination's supported constraints. BioSeqDB table
restrictions still apply, including unsupported foreign keys and
\texttt{AFTER ROW} triggers. This is distinct from the specialized
\texttt{bioseq.load} bulk path discussed in the local-ingestion section.
Once imported, subsequent searches read the local sequence.

Local \texttt{mmsearch} accepts prepared parameters, including generic plans.
Pass a sequence with an explicit alphabet and storage unit. The example query
below is the rehearsal query. Replace it with the biological query of interest
when searching the imported region. A live region need not contain this motif.

\par\noindent\begin{minipage}{\linewidth}
\begin{lstlisting}[style=bioseq-appendix-sql]
PREPARE search_imported(bioseq.seq) AS
WITH local_matches AS MATERIALIZED (
  SELECT provider, accession, start0,
         bioseq.target_start(accession) AS match_start_in_region,
         bioseq.target_end(accession) AS match_end_in_region,
         bioseq.cigar(accession) AS cigar
  FROM imported_provider_targets
  WHERE public.mmsearch(sequence, $1, aligner_options(best_n => 3))
)
SELECT provider, accession, match_start_in_region,
       match_end_in_region, start0 + match_start_in_region
         AS provider_start0,
       start0 + match_end_in_region AS provider_end0, cigar
FROM local_matches
ORDER BY provider, accession, provider_start0;

EXECUTE search_imported(
  'TAGTCACTGTAGCCTGATCGTAGCTAGTCA'::bioseq.seq('dna4', 64)
);
DEALLOCATE search_imported;
\end{lstlisting}
\end{minipage}

Adding \texttt{start0} is correct for a stored forward-oriented region. If an
application stores the reverse complement, it must retain that orientation
and transform match endpoints accordingly. A chunk boundary also cuts a
sequence into separate targets unless the application reconstructs the
record or deliberately overlaps chunks. Independent chunk searches can miss
matches spanning a boundary.

\subsection{Perform bulk maintenance and retain provenance}

For many records, import into a fresh staging table and perform one explicit
minimizer-index update after the batch. The following complete example creates
a new generation from the local rows already imported. Replace its
\texttt{INSERT SELECT} with bounded provider imports for a fresh remote
generation. \texttt{SET LOCAL} restores the maintenance setting when the
transaction ends.

\par\noindent\begin{minipage}{\linewidth}
\begin{lstlisting}[style=bioseq-appendix-sql]
BEGIN;
CREATE TABLE imported_provider_next
  (LIKE imported_provider_targets INCLUDING ALL) USING bioseq;
SET LOCAL bioseq.index.insert_maintenance.enabled = off;

INSERT INTO imported_provider_next
SELECT * FROM imported_provider_targets;

SELECT count(*) AS rows, sum(bioseq.len(sequence)) AS bases
FROM imported_provider_next;

SELECT bioseq.update_index(
  'imported_provider_next'::regclass, 'minimizers_per_fragment'
);
COMMIT;
\end{lstlisting}
\end{minipage}

Before publishing a generation, compare accession coverage, row counts, and
lengths with the requested input, apply the application's validation policy,
and complete the index update. Then switch the application's selected
generation in a separate controlled step. Creating a staging table does not
automatically publish it to readers.

Store provider metadata with the sequence when its origin matters. The next
alternative imports directly into a table that retains the returned version,
provenance JSON, and retrieval timestamp.

\par\noindent\begin{minipage}{\linewidth}
\begin{lstlisting}[style=bioseq-appendix-sql]
CREATE TABLE imported_provider_provenance (
  provider text NOT NULL,
  accession text NOT NULL,
  start0 bigint NOT NULL,
  end0 bigint NOT NULL,
  sequence bioseq.seq('dna4', 64) NOT NULL,
  provider_version text,
  provenance jsonb NOT NULL,
  retrieved_at timestamptz NOT NULL,
  PRIMARY KEY (provider, accession, start0, end0)
) USING bioseq;

INSERT INTO imported_provider_provenance
  (provider, accession, start0, end0, sequence,
   provider_version, provenance, retrieved_at)
SELECT provider, accession, start0, end0,
       sequence::bioseq.seq('dna4', 64),
       provider_version, provenance, retrieved_at
FROM bioseq_external.sequences
WHERE provider = 'ucsc'
  AND assembly_database = 'hg38'
  AND chromosome = 'chrM'
  AND start0 = 0 AND end0 = 100
  AND sequence::text ~ '^[ACGT]*$';
\end{lstlisting}
\end{minipage}

The adapter's provenance may contain only the provider name, and the version
may be null. Retain the request coordinates and assembly in application
metadata as well. A retrieval timestamp records when data was received. It
does not establish a consistent snapshot across several remote requests.

\subsection{Protein limits, failure handling, and cleanup}

Protein discovery remains useful even though protein retrieval and a native
protein alignment workflow are unavailable through these sequence relations.
For example, inspect the metadata and review status of a UniProt record:

\par\noindent\begin{minipage}{\linewidth}
\begin{lstlisting}[style=bioseq-appendix-sql]
SELECT accession, title, sequence_length, is_reviewed, retrievable
FROM bioseq_external.catalog
WHERE provider = 'uniprot' AND accession = 'P04637';
\end{lstlisting}
\end{minipage}

\begin{itemize}
  \item For remote failures, check the requested operation, identifiers,
        provider availability, rate limits, and cancellation. Follow provider
        usage policies: the SQL path does not supply a shared rate controller
        or a general remote-result cache. NCBI E-utilities and UCSC publish
        their own request guidance~\autocite{ncbiEutilitiesUsage,ucsc_rest_api}.
  \item The direct NCBI stream and archive routes retry selected transient
        failures up to five attempts, only before any sequence row has been
        emitted. A failure after emission is terminal for that request.
        Cancelled or failed application imports should be retried as a fresh
        staging generation, with coverage checked before publication.
  \item For local search failures, inspect the local table access method,
        target and query alphabets, index-maintenance state, and search plan.
        A provider outage cannot change an already committed local copy.
  \item Record requested accessions, local import counts, validation results,
        index completion, and the local search plan. Keep remote request
        diagnostics separate from local query measurements. Omit secrets and
        provider continuation tokens from shared logs.
\end{itemize}

The examples explicitly deallocate their prepared statements. After finishing
the scratch workflow, remove its local tables and restore the session settings.
This cleanup does not remove the installed external relations or application
role grants.

\par\noindent\begin{minipage}{\linewidth}
\begin{lstlisting}[style=bioseq-appendix-sql]
DROP TABLE IF EXISTS imported_provider_next;
DROP TABLE IF EXISTS imported_provider_provenance;
DROP TABLE IF EXISTS imported_provider_targets;
DROP TABLE IF EXISTS imported_ncbi_records;
RESET bioseq.index.insert_maintenance.enabled;
RESET bioseq.vectorization.enabled;
\end{lstlisting}
\end{minipage}
